\customchapter{ABSTRACT} 


Stone heritage is exposed to microbial biodeterioration, a process by which the communities colonizing the rock dissolve its mineral matrix and alter its surface. The Raimondi Stela, a granite lithosculpture of the Chavín culture held in a national museum, exhibits biofilms and chromatic alterations whose prior assessment was limited to taxonomic identification, without clarifying the deterioration capacity of the living organisms inhabiting it. The objective of this project was to design a computational analysis workflow to characterize the biodeterioration potential and the genomic adaptation signals of the dominant bacterial lineages of the monument's viable fraction, in order to support risk assessment and genomics-based preventive conservation. High-throughput metagenomic sequencing data obtained from the cultivable surface fraction were used as the starting point, upon which a reproducible workflow was structured, integrating genomic quality control, whole-genome-based taxonomic classification, the grouping of redundant genomes into lineages, comparative pangenomic reconstruction against global references, and functional annotation of metabolic potential.  \\


\noindent \textbf{Keywords:}\\
\noindent  Biodeterioration; Stone heritage; Metagenomics; Pangenome;
Metabolic potential; Preventive conservation; Raimondi Stela


